\section{Trace presentations after a change of middle interval}

\begin{theorem}[Trace presentation of one transferred middle factor]
    \label{thm:interval_trace_factor_transfer}
    \lean{TNLean.MPS.AreaLaw.exists_traceChainOperator_changingRegionalMiddleFactor}
    \leanok
    \uses{thm:interval_transfer_entry_coordinates,
        thm:trace_interval_slice_configuration_extend}
    Let $q\ge1$ and let $E=\lbrack c,d)$ and $J=\lbrack u,v)$ satisfy
    (\ref{eq:interval_mpo_transfer_endpoints}). Supply a matrix $V$ on
    $\mathcal C(E)$ with a trace presentation of virtual dimension $\chi$.
    Fix any row and column configurations on $E\setminus J$ and
    $J\setminus E$, and let $W$ be the resulting prescribed regional
    middle factor. There are tensors on $J$ of the same dimension $\chi$
    and a closing matrix whose trace contraction equals $W$.
    The tensors and closing matrix are chosen for the whole operator,
    before its variable row and column configurations. The dimension
    $\chi$ and the overlap length may vanish.
\end{theorem}
\begin{proof}\leanok
    \uses{thm:interval_transfer_entry_coordinates,
        thm:trace_interval_slice_configuration_extend}
    Write $W$ in the natural conditional form
    (\ref{eq:interval_transfer_entry_coordinates}). Substitute the trace
    presentation of $V$ and apply
    Theorem~\ref{thm:trace_interval_slice_configuration_extend}.
\end{proof}

\begin{theorem}[A common transferred family at the same trace dimensions]
    \label{thm:interval_trace_common_family_transfer}
    \lean{TNLean.MPS.AreaLaw.changingRegionalMiddleFamily_trace_and_global_decomposition}
    \leanok
    \uses{thm:interval_trace_factor_transfer,thm:regional_middle_exterior_transfer}
    Under (\ref{eq:interval_mpo_transfer_endpoints}), let $q\ge1$ and let
    $V_t$, $t\in\iota$, be a finite family on $\mathcal C(E)$, each with
    a trace presentation of virtual dimension $\chi_t$. Define the same
    prescribed carrier $I$ and matrices $W_r$ by
    (\ref{eq:interval_mpo_transfer_index}) and
    (\ref{eq:interval_mpo_transfer_factor}). Then
    $|I|=|\iota|q^{2|E\setminus J|+2|J\setminus E|}$, and the member
    indexed by $(t,\alpha_r,\alpha_c,\eta_r,\eta_c)$ has a trace
    presentation of dimension exactly $\chi_t$ on $J$.
    The family is chosen before every exterior family. For every later
    exterior family $M_t$, there is an exterior family $M'_r$ giving the
    same whole-chain operator as in
    (\ref{eq:interval_mpo_transfer_decomposition}). Empty carriers and
    empty overlaps are allowed.
\end{theorem}
\begin{proof}\leanok
    \uses{thm:interval_trace_factor_transfer,thm:regional_middle_exterior_transfer}
    Apply Theorem~\ref{thm:interval_trace_factor_transfer} to each prescribed
    member. The two configuration pairs give the stated cardinality.
    The regional exterior-transfer identity gives the decomposition for
    every later exterior family, without changing the middle matrices.
\end{proof}
